% Mobile deployment – artifacts and sizes first

% Auto-generated mobile artifact table (if present)
% Auto-generated mobile artifact sizes
\begin{table}[h]
  \centering
  \caption{Exported on-device model artifacts.}\label{tab:mobile_artifacts_auto}
  \begin{tabular}{lcc}\toprule
  Artifact & Format & Size (MB) \\ \midrule
  Stage~1 (MobileNetV4) & ONNX & 37.30 \\ 
  Stage~2 (EfficientNetV2-B3) & ONNX & 49.55 \\ 
  \textbf{Total} & -- & \textbf{86.85} \\ 
  \bottomrule\end{tabular}\end{table}

% Auto-generated on-device detection summary
\begin{table}[h]
  \centering
  \caption{On-device cascade runtime and usage (auto).}\label{tab:mobile_runtime_auto}
  \begin{tabular}{lrrrrrr}\toprule
  Mode & $N$ & Deepfake Rate & Stage2 Rate & Preproc (ms) & Inference (ms) & Total (ms) \\ \midrule
  Images (single/batch) & 152 & 0.572 & 0.072 & 3.4 & 176.3 & 179.7 \\ 
  \bottomrule\end{tabular}\end{table}

% Auto-generated on-device accuracy summary (requires folder-structured datasets)
\begin{table}[h]
  \centering
  \caption{On-device accuracy and error rates (auto).}\label{tab:mobile_accuracy_auto}
  \begin{tabular}{lrrrrrrr}\toprule
  Mode & $N_{lbl}$ & Acc & FNR & FPR & Prec & Rec & Stage2 Rate \\ \midrule
  Images (single/batch) & 152 & 0.928 & 0.025 & 0.125 & 0.897 & 0.975 & 0.072 \\ 
  \bottomrule\end{tabular}\end{table}


\paragraph{PTQ vs. QAT.} We adopt post\hyp training dynamic quantization (PTQ) for linear layers due to its simplicity and strong accuracy\hyp size trade\hyp off on our models. Quantization\hyp aware training (QAT) \cite{jacob2017integeronly} can further reduce the quantization gap by learning with simulated quantization noise, at the cost of training complexity and longer cycles; see \cite{gholami2021quant_survey,han2016deepcompression,esser2020lsq} for broader comparisons of PTQ/QAT. As complementary levers, pruning and knowledge distillation reduce compute and model size with modest accuracy loss \cite{blalock2020pruningsurvey,hinton2015distilling,romero2015fitnets,furlanello2018bornagain}. We plan to evaluate per\hyp channel INT8 and selective INT4 for heads/backbone under the same export pipeline.

% Placeholders for on-device latency table (to be filled after measurement)
\begin{table}[h]
  \centering
  \caption{On-device performance (to be completed): average per-frame latency and Stage~2 escalation rate.}
  \label{tab:mobile_perf}
  \begin{tabular}{lcccc}
    \toprule
    Device & Stage~1 (ms) & Stage~2 (ms) & S2 Rate (\%) & APK Size (MB) \\
    \midrule
    Mid-range & -- & -- & -- & -- \\
    High-end  & -- & -- & -- & -- \\
    \bottomrule
  \end{tabular}
\end{table}
